\section{Every spin alternating layer with an entire tensor partner}
\label{sec:spin-alternating-layers}

\paragraph{Status, authorship, and precise source inputs.}
HD-FT23--26, reserved question Q2dz, is a new derivation by
\texttt{/root/g34\_eightfolds}.  The block restriction argument below
applies to every faithful spin irreducible, not just the basic spin
characters.  Its source inputs are the following published results.
Guralnick--Tiep, Lemma~2.23, with its complete proof on pp.~621--622,
gives the character-ratio bound $7/8$ for every faithful irreducible
of $2.A_m$ or $2.S_m$, $m\ge6$.
Their Proposition~3.4 proof, p.~626, explicitly records the complex
spin minimum $2^{\lfloor m/2\rfloor-1}$.
\cite{HD:GT2012}
The underlying Kleshchev--Tiep Theorem~A, its degree definition (1.1),
and the complete Section~3 proof have also been read in the original
author version, numbered pp.~2 and 9--12.  Its opening Schur-cover
presentation, numbered pp.~1--2, fixes the conventions used here.
\cite{HD:KT2004}
These are inputs, not new proofs of the minimum-degree theorem or its
upstream modular branching results.  The column-orthogonality identity
is Serre, Proposition~7, p.~20, whose complete proof has been read.
\cite{OG:Serre1977}
The source/status record distinguishes the actual published and author
versions and the unsuccessful publisher-PDF access attempt.  The
complete new proof is written below. The author, root and independent
whole-proof, source and code readings pass without mathematical repair;
their distinct scopes and exact file bindings are recorded separately.

We use the double cover $\widehat S_m$ in which adjacent transpositions
lift to involutions $t_i$, with central involution $z$ and relations
\[
 t_i^2=1,\quad t_it_{i+1}t_i=t_{i+1}t_it_{i+1},\quad
 t_it_j=z t_jt_i\quad (|i-j|>1).
 \tag{SPA.1}
\]
Write $\widehat A_m$ for the inverse image of $A_m$.
The other Schur-cover convention is equivalent on the spin algebra
over $\mathbf C$ by multiplying the generators by $i$; only scalar
phases change, so absolute traces and projective support do not.
The source bounds therefore apply with (SPA.1).
Throughout, a spin module means that $z$ acts as $-I$, and
$\nu(T)=\dim T-\max_\alpha\dim\ker(T-\alpha I)$.

\subsection*{HD-FT23. An all-spin exponential bound from disjoint cycle blocks}

\begin{lemma}[HD-FT23]
Let $n\ge8$, let $V$ be any nonzero complex spin module of
$\widehat A_n$ or $\widehat S_n$, with character $\chi$ and degree $a$,
and let $g$ project to $x\in S_n$.  Put $s=n-f(x)$, the number of
moved letters, and $r=7/8$.  Then
\[
 \frac{|\chi(g)|}{a}\le r^{(s-13)/10}.
 \qquad
 x\ne1\quad\Longrightarrow\quad\frac{|\chi(g)|}{a}\le r.
 \tag{SPA.2}
\]
No irreducibility of the restriction to a cycle-block subgroup is
assumed.  In particular this estimate includes every faithful spin
irreducible, every split spin constituent, and every spin extension
to the chosen double cover of $S_n$.
\end{lemma}

\begin{proof}
All modules under consideration may be made unitary.  A spin
irreducible of $\widehat A_m$, $m\ge6$, is faithful: its kernel is
normal, its projection to the simple quotient is either trivial or
surjective, and the latter would give a nontrivial abelian quotient
of the perfect double cover while killing no $z$.  For
$\widehat S_m$, the kernel centralizes its faithfully represented
$\widehat A_m$, hence projects into the trivial centralizer of
$A_m$ in $S_m$ and again is trivial.  Thus the published $7/8$
bound applies to every spin constituent, and by a convex combination
it applies to every spin module.  This proves the second assertion.

We first bound the ratios on the blocks to be used.  Every such
block has size $m\ge6$ and its nonidentity permutation is even,
except for at most one single even-length cycle.  The relevant
group is $\widehat A_m$ for an even block and $\widehat S_m$ for
that possible odd block.  For $6\le m\le10$, the ratio is at most
$r\le r^{m/10}$ by the preceding published bound.

For the larger blocks the permutation is either a single cycle
or a pair of even-length cycles.  Column orthogonality and the
spin minimum degree give, for each faithful irreducible $\psi$,
\[
 |\psi(h)|\le\sqrt{|C_{\widehat H_m}(h)|},\qquad
 \psi(1)\ge2^{\lfloor m/2\rfloor-1}\quad(m\ge8).
 \tag{SPA.3}
\]
The centralizer of a lift maps into the permutation centralizer
with kernel of size at most two.  For a single $m$-cycle this
gives $|C_{\widehat H_m}(h)|\le2m$.  For two even cycles of
lengths $u,v$, $u+v=m$, it gives at most $2uv$ when $u\ne v$
and at most $4uv$ when $u=v$; in either case the result is at
most $m^2$.  The possible factor two from the double cover is
included in these bounds.

If $m\ge10$ is even, both a single cycle and a pair of even
cycles therefore have ratio at most $2m\,2^{-m/2}$.
At $m=10$ this is $5/8<r$.  On increasing $m$ by two, the
ratio of this bound divided by $r^{m/10}$ changes by the factor
\[
 \frac{1+2/m}{2r^{1/5}}\le\frac{24}{35}<1.
 \tag{SPA.4}
\]
Here $r^{1/5}\ge r$ and $1+2/m\le6/5$.
For odd $m\ge11$, the block is a single odd cycle and the
ratio is at most $4\sqrt m\,2^{-m/2}$.  Its value at $m=11$
is less than $1/2$, since its square is $11/128<1/4$.
Also $1/2<r^2<r^{11/10}$.  The corresponding step by two
has squared factor at most
\[
 \frac{13/11}{4r^{2/5}}\le\frac{208}{539}<1.
 \tag{SPA.5}
\]
Consequently every larger block also has ratio at most $r^{m/10}$.
All these inequalities apply to reducible spin restrictions by
the same convexity argument.

Here is a cycle-block construction covering all permutations.
Ignore fixed letters.  Treat each odd cycle of length at least
seven as one block.  Pair the cycles of lengths three and five;
each resulting block has size six, eight or ten and is even.
At most one such short odd cycle is left aside, on at most five
letters.

Among the even-length cycles, if their number is odd, first take
one cycle out.  If its length is at least six it is the single
possible odd block; if its length is two or four leave it aside.
Pair all remaining even-length cycles.  Each pair has even size
at least six except a pair of two transpositions, of size four.
Merge two size-four pairs into a size-eight even block, and leave
at most one size-four pair aside.  Thus every charged block has
size at least six, there is at most one odd block, and the total
number of letters left aside is at most
\[
                         5+4+4=13.
 \tag{SPA.6}
\]
For $x\in A_n$ there is no odd block at all.  The merged
size-eight block, which has four transpositions, uses only the
$7/8$ bound; it is not put into the two-cycle centralizer estimate.

The lifted block groups commute.  To verify this directly,
conjugate disjoint blocks to consecutive intervals.  Their internal
adjacent generators have index difference at least two, so (SPA.1)
shows that lifts of disjoint permutations commute with sign
$z^{\epsilon_1\epsilon_2}$, where $\epsilon_i$ is their parity.
Every charged group is the even group $\widehat A_m$ except
possibly one $\widehat S_m$.  Thus all charged groups commute.
Their common central $z$ identifies them as a quotient of their
direct product by central relations.  The residual permutation
on the uncharged letters commutes with each charged group as
well: if an odd charged block exists, the residual permutation
is even; otherwise every charged group is even.

Write $g$ as a product of lifts of its charged block permutations
and a residual lift, up to the harmless central sign.  Restrict
$V$ to the product of the charged groups.  Complete reducibility
and the irreducible direct-product theorem decompose it into
isotypic spaces
\[
                   \bigotimes_j V_j\otimes M.
 \tag{SPA.7}
\]
Every $V_j$ is spin, since its own $z$ is the same global
involution acting as $-I$.  The residual lift commutes with the
whole product, so on each such isotypic space it acts as the
identity on the irreducible tensor product and as a unitary
operator on $M$.  Its trace has absolute value at most $\dim M$.
The product of the charged-block ratios is therefore at most
$r^{\sum_j m_j/10}$ on each summand.  Since
$\sum_jm_j\ge s-13$, summing their traces proves the first
inequality in (SPA.2).  Empty charged products cause no problem:
then $s\le13$ and the asserted upper bound is at least one.
\end{proof}

\subsection*{HD-FT24. Compatible outer operators and an explicit cubic bound}

\begin{lemma}[HD-FT24]
Let $n\ge65536$ and let $\rho:\widehat A_n\to\operatorname{GL}(W)$
be a faithful irreducible spin representation of degree $a$.
Every compatible projective normalizer operator, after removal of
a scalar, is a finite-order operator $T$ with an actual natural
permutation $x\in S_n$ such that, for every integer $j$,
\[
 \frac{|\operatorname{tr}(T^j)|}{a}
 \le \left(\frac{\max\{f(x^j),1\}}n\right)^3,
 \qquad a\ge n^3.
 \tag{SPA.8}
\]
For $x\ne1$ this implies the all-target local estimate
\[
 \nu(T)\ge A\operatorname{ind}(x)+B,
 \qquad
 A=\frac{a(n-4)}{n^2},\quad B=\frac an.
 \tag{SPA.9}
\]
The natural permutation is supplied by conjugation on the simple
core quotient, not by a branchwise choice of a projective lift.
\end{lemma}

\begin{proof}
First establish (SPA.8) for any spin character of either double
cover.  For a nonidentity permutation put $s=n-f$.  If $s\le n/24$,
then
\[
 \frac78<\left(\frac{23}{24}\right)^3,
 \qquad 23^3-\frac78\,24^3=71,
 \tag{SPA.10}
\]
so the second bound in (SPA.2) is sufficient.  If $s\ge n/24$,
the first bound, together with $\log(7/8)\le-1/8$, gives
\[
 \frac{|\chi(g)|}{\chi(1)}
 \le\exp\left(-\frac n{1920}+\frac{13}{80}\right)
 <n^{-3}\le\left(\frac{\max\{f,1\}}n\right)^3.
 \tag{SPA.11}
\]
Here is the full endpoint calculation.  The first four terms of
the exponential series give $\exp(7/10)>2$, hence $\log2<7/10$.
At $N=65536=2^{16}$,
\[
 \frac N{1920}-\frac{13}{80}-3\log N
 >\frac{512}{15}-\frac{13}{80}-\frac{168}{5}
 =\frac{89}{240}>0.
 \tag{SPA.12}
\]
The derivative of $n/1920-13/80-3\log n$ is positive for
$n>5760$, so this holds for every $n\ge N$.
An identity permutation has a central lift and ratio one,
which is exactly the right side of (SPA.8).
The minimum in (SPA.3) gives $a\ge n^3$ throughout this range.
For completeness this elementary comparison already holds for
all integers $n\ge64$: check $64$ and $65$, then increase $n$
by two, when the power of two doubles and
$((n+2)/n)^3\le(33/32)^3<2$.

We next check all compatible outer operators.  Projective
normalization of the represented perfect group is linear
normalization, by taking derived groups and applying the inclusion
also to the inverse operator.  An automorphism of $\widehat A_n$
acting trivially on $A_n$ differs from the identity by a homomorphism
$\widehat A_n\to\langle z\rangle$, hence is trivial by perfectness.
Together with the accepted $\operatorname{Aut}(A_n)=S_n$ input,
this identifies every compatible automorphism with conjugation
by a lift of its actual permutation $x\in S_n$.

If no odd automorphism stabilizes $\rho$, every compatible operator
is a scalar times an inner matrix.  Otherwise compose one odd
intertwiner with an inner matrix to obtain an intertwiner $S$ for
the chosen transposition lift $t_1$.  Schur's lemma makes $S^2$
scalar; rescale $S$ so that $S^2=I=\rho(t_1^2)$.  The conjugation
relations and this square relation give a genuine representation
of the index-two extension $\widehat S_n$ extending $\rho$.
It is irreducible because its restriction is, and faithful because
any kernel would centralize the faithfully represented core and
then lie in $\langle z\rangle$.
Every compatible odd operator now differs from its representing
matrix by a scalar, by Schur's lemma.  Thus in both cases it has
a representative in one fixed finite double-cover representation.
Every power projects to $x^j$, so the proved absolute character
bound gives (SPA.8) for all its powers, including central powers.
No formula limited to basic spin or one of the two outer cosets
has been substituted for this extension argument.

We repeat the short all-target Fourier argument to make its
projective scope explicit.  Let $I=\operatorname{ind}(x)>0$.
Fourier averaging over the actual finite order of $T$ bounds each
target multiplicity by the average of $|\operatorname{tr}(T^j)|$.
The average number of fixed natural letters of $x^j$ is $n-I$,
even if $T$ has twice the order of $x$.
If $I\ge n/3$, use
$(\max\{f,1\}/n)^3\le f/n+n^{-3}$; after subtracting $AI/a$
the support margin is at least $4I/n^2-n^{-3}>1/n$.
If $3\le I<n/3$, every power fixes at least $n-2I>n/3$
letters, so $1-u^3\ge(13/9)(1-u)$ and the margin is at least
$4I/(9n)+4I/n^2>1/n$.
Finally for $I=1,2$, every power has $u\ge1-4/n\ge5/6$,
and $1+u+u^2\ge91/36>5/2$; the margin is at least
$3I/(2n)+4I/n^2>1/n$.  These cases prove (SPA.9) for every
eigenvalue target.  Scalar changes do not change $\nu$.
\end{proof}

\subsection*{HD-FT25. Every tensor count with the whole external partner}

\begin{theorem}[HD-FT25]
Let $n\ge65536$.  Suppose an actual finite coefficient image has
a stable whole spin layer with transitive component permutation,
and an actual entire external partner, on
\[
 V=\left(\bigotimes_{j=1}^t W_j\right)\otimes U,
 \quad t\ge1,\quad \dim W_j=a,\quad b=\dim U,\quad e=a^tb.
 \tag{SPA.13}
\]
The represented layer is an irreducible normal central product of
$t$ faithful double-cover components $\widehat A_n$; normality
here refers to its actual projective layer image.  Suppose every
nonscalar operator in the entire partner image satisfies
\[
                         \nu_U(S)/b\ge1/n.
 \tag{SPA.14}
\]
For $b=1$ this is vacuous.  There is no connected sphere cover
with $r\ge12$ genuine coefficient inertias and moving-rank sum
$D=2e+6$ in this setting.  In particular the layer cannot occur
in an actual rank-nine sixfold supplier.  All original fields,
all compatible outer operators and all component permutations
are allowed.
\end{theorem}

\begin{proof}
Transitivity implies that the factor representations belong to
one automorphism orbit.  Choose fixed component identifications,
using odd identifications where necessary, so that every $W_j$
is one chosen $W$ up to such a fixed gauge.  This only conjugates
the natural block action.  Each compatible cycle product is now
a normalizer operator on $W$ and has the finite representative
and all-target bound of HD-FT24.

Here are the actual quotient homomorphisms.  Conjugation of the
represented central product gives its component permutation and
its automorphisms on the simple quotients.  They give a genuine
homomorphism from the actual coefficient image into
$S_n\wr S_t$.  Its kernel on the layer is scalar: a component-
trivial core automorphism differs by a central character, which
perfectness kills, and Schur's lemma kills its remaining linear
centralizer.  The partner does not change this homomorphism.
Consequently the actions on $t$ factors and $nt$ natural letters
are actual quotient actions.  The latter is transitive because
the layer contains $A_n^t$, transitive inside every block, and
is transitive on blocks.  No linear representation of $S_n$
on the spin module and no individual stable factor quotient
for $t>1$ is assumed.

At a factor-fixing entry let $I$ be the sum of the active
within-factor indices.  If there are $m>0$ active factors,
spectral convolution and the average of their normalized local
bounds give at least $AI/(ma)+B/a\ge AI/(ta)+B/a$.
If all projective core factors are identity, a nonscalar partner
pays $1/n=B/a$; if it is scalar too, genuineness gives full
actual moving rank.  Thus every genuine factor-fixing entry pays
\[
             \operatorname{rank}(1-T)/e\ge AI/(ta)+B/a.
 \tag{SPA.15}
\]

For $t=1$, the actual connected $n$-letter quotient has total
index at least $2n-2$.  Summing (SPA.15) gives directly
\[
 D-2e\ge\frac ea\{2A(n-1)+12B-2a\}
       =\frac{e(2n+8)}{n^2}>6,
 \tag{SPA.16}
\]
since $e\ge a\ge n^3$.  Hence the single stable spin factor
is included, even with an arbitrary covered partner.

Now let $t\ge2$.  The accepted weighted-word argument gives
every factor-moving entry support at least $(a-1)/(2a)\ge B/a$.
Four moving entries and eight further genuine entries would give
$D\ge2e+(8B-2)e/a>2e+6$, because $e/a\ge a$ and
$a(8B-2)>6$.  There are at most three movers and at least nine
factor-fixing entries.

For a general entry with factor permutation $\pi$, cycle lengths
$\ell_j$ and compatible cycle products $x_j$, its actual
$nt$-letter index and weighted-cycle bounds are
\[
 \begin{gathered}
 I=n\operatorname{ind}(\pi)+\sum_j\operatorname{ind}(x_j),\\
 s_j\ge1-\frac1{\ell_j}
       +\frac{\nu(T_{x_j})}{\ell_j a}
       -\frac{R_{\ell_j}(a)}{\ell_j a^{\ell_j}},\qquad
 R_1=0,\quad R_\ell(a)/a^\ell\le1/a\ (\ell\ge2).
 \end{gathered}
 \tag{SPA.17}
\]
The exact $R_\ell=\sum_{j=1}^{\ell-1}a^{\gcd(\ell,j)}$
is the full D-ATT08/HD-TW01 defect; no short-orbit factor is
omitted.  Its hypotheses require only the actual diagonalizable
cycle product and $a\ge4$, so scalar projective lifts cause
no change.  Convolution bounds full support below by the weighted
average $\sum_j(\ell_j/t)s_j$.  Using the weaker
$\nu(T_{x_j})\ge A\operatorname{ind}(x_j)$, valid also for
identity products, and counting at most $\operatorname{ind}(\pi)$
nontrivial cycles yields
\[
 \frac{\operatorname{rank}(1-T)}e
 \ge\frac{AI}{ta}+\frac\lambda t\operatorname{ind}(\pi),
 \quad\lambda=\frac4n-\frac1a>0,
 \quad L=a-(n-1)A=\frac{a(5n-4)}{n^2}\ge1.
 \tag{SPA.18}
\]
The $B/a$ term in (SPA.15) is retained at the factor-fixing
entries only.  The actual connected quotient covers give total
indices at least $2nt-2$ and $2t-2$.  Summation therefore gives
\[
 \begin{split}
 D-2e&\ge\frac ea\left(r_{\rm fix}B-2-\frac2t(L-1)\right)\\
     &\ge\frac ea(9B-L-1)
       =\frac ea\left(\frac{4a(n+1)}{n^2}-1\right)>6.
 \end{split}
 \tag{SPA.19}
\]
The last inequality follows from $e/a\ge a\ge n^3$.
This is the required contradiction for every $t\ge2$.

Finally the accepted actual intermediate-cover interface applies
because $e>6$: base genus at least two contradicts the multiplicity
Euler identity, and genus one has only six moving-rank units and
cannot give period rank nine.  On the sphere that rank requires
at least twelve genuine coefficient marks and the exact budget
$D=2e+6$.  These are necessary geometric facts, not a construction
of a source from numerical data.
\end{proof}

\subsection*{HD-FT26. The maximal spin layer in a complete primitive decomposition}

\begin{corollary}[HD-FT26]
In the complete recursive decomposition D-TD12 of an actual finite
primitive coefficient image, suppose one extracted layer is spin
alternating of natural degree $n\ge65536$, and that $n$ is maximal
among all alternating natural degrees in every recursion layer.
Then that layer cannot occur in an actual rank-nine sixfold curve
supplier, for any tensor count and any original field.
\end{corollary}

\begin{proof}
Take the entire transitive layer as the first factor of (SPA.13).
All remaining extracted pieces form its whole partner, with no
unaccounted multiplicity space; their full matrix algebras remain
normalized or permuted by the actual image.  D-TD12 supplies
genuine finite projective quotients at every extraction step, so
the conjugation actions used in HD-FT25 pull back to the original
coefficient image.  The complete partner theorem HD-FT21 gives
\[
 \nu_U(S)/\dim U\ge
 \frac1{\max\{80,m_{*,U}-1\}}\ge\frac1n
 \tag{SPA.20}
\]
for each nonscalar partner operator, by the maximality hypothesis.
For an empty partner the assertion is vacuous.  HD-FT25 applies.
\end{proof}

\paragraph{Remaining scope and exceptional covers.}
This excludes a spin layer when it has the maximal alternating
natural degree in the exhibited complete primitive decomposition;
without maximality the explicit whole-partner hypothesis (SPA.14)
is retained.  It does not assert that every individual component
of a permuted layer is stable.  The exceptional Schur covers at
$A_6$ and $A_7$ remain the separately treated finite families.
Their natural double-cover subgroups in the block proof use only
the exact $m\ge6$ scope of the published $7/8$ theorem.  No
assertion about the exceptional triple or sixfold covers is
deduced from that block argument.  No original induction-index,
field-descent, polarization, or existence conclusion is inferred.
